\documentclass{article}
\usepackage{amsmath}
\begin{document}
\title{A qualified historical assignment does not transport the replacement body's realized motion}
\maketitle

\paragraph{Claim and scope.}
In a pinned Robot Warehouse (RWARE) fixture, two histories give the successor the
same old flag, the same authentic as-of dispatch, the same commissioned mirror
placement, the same issued actions on both the successor and a nearby body, and
the same complete normalized-dictionary observation at tick 21. Nevertheless,
under a common two-action deadline, exactly one of two predeclared inspection
targets is available in each history, and the available target differs. The
one-bit settled outcome of the successor's named step distinguishes the two.
The binding of the bit to the actual body and computation cut is a separate
requirement. No claim of a field fault, general transport, or I1/I2 acceptance.

\paragraph{Fixture and cut.}
Continue the exact native two-delivery source in the parent replay. The first
delivery at tick 1 sets the shelf assignment; the original body at tick 2 sees
the corresponding mark, then returns to row 3. The second actual delivery is
complete at tick 15 and both worlds now request item 4. The standby begins at
the reflected column (map $x\mapsto6-x$), row 3 and heading DOWN; its commissioning
relation is trusted in this experiment. The first freight carrier turns twice
and leaves the visible area. At step 19 the standby requests FORWARD in both
worlds. A stationary body at row 4 cancels that request in one world; the
same body starts one cell farther south in the other and does not cancel it.
The other body subsequently requests two forwards and is at row 6 in both
worlds by tick 21. The standby was not reading during steps 19--21. The complete
normalized observation at 21 agrees (including native heading and local patch),
and the earlier marked shelf is out of view. Neither the other body's full
pre-move state nor a continuously available out-of-band encoder is part of the
input. The fresh absolute-output control does provide the latter distinction
from the same pin and simulator state.

\paragraph{Exact test.}
At tick 21 the authority chooses at most one of the two predeclared requests,
which must reach $(3,3)$ or $(3,4)$ by tick 23, respectively. In each test the
goal is supplied, not opportunistically redefined. From either initial column,
the matching row needs one quarter turn and a forward. From row 3, reaching
$(3,4)$ requires at least two translations and a turn. From row 4 facing DOWN,
reaching $(3,3)$ requires a north translation, a horizontal translation, and at
least two turns. These bounds are attained. The table is
\[
\begin{array}{c|cc}
 & (3,3)&(3,4)\\\hline
\mathrm{canceled}&2&3\\
\mathrm{effective}&4&2
\end{array}
\]
in native steps. The settled row is thus the exact quotient for these two
feasibility queries once the original as-of evidence, commissioning map,
heading and static suffix context are supplied: constant messages cannot
answer both histories, a single Boolean can. This is not an absolute lower bound
on pose-estimator memory, and a resource-aware direct controller given a fresh
absolute observation, or the same qualified outcome, ties. The comparison with
using only the request is weaker: it declines the feasible row-3 request when
the movement was canceled and admits the infeasible row-4 request in that case.
Pretending freshness for the stationary cut produces the symmetric error.
A correctly labeled old cut, a different body, or a different appointment is
not accepted as the current body record. Native enumeration checks all 340
suffixes of lengths 1--4 in each of eight source worlds. Unrelated native
delivery succeeds before either branch. All suffix checks are one-mover,
empty-handed, within a fixed obstacle map over the two-step commitment;
a movement after this cut, a new appointment, an unreported obstruction, or
a broken mirror calibration requires new evidence. Three successive cuts
are captured from real native before/after transitions, with no hardware claim.

\paragraph{Radio join.}
The second replay sends the historical assignment (five bytes) from P, the
independent actual second-delivery report (three bytes, due at 15) from Q, and
the one-byte settled outcome from E, released downstream of the native
step and due before admission at tick 21. A preinstalled source namespace at
the broker is sufficient for the test: E is bound to the one appointed body,
named cut and effect kind. The record has no self-describing identity,
authentication, retransmission costs, generalized concurrent tags or epoch-rollover
guarantee. The initial historical dispatch is available from tick 1;
P is a transmission of that old record, not the latest dispatch.
Both control policies pay the same receive budget, and the receiver uses actual
source-tagged application callbacks for its decisions, not an assumed successful
transmission. The time mapping is tick $t/16$ seconds, a selected high-signal
SubT broker/CommsClient path with the pinned rolling channel (including age
exactly 1 s), rather than rendered RF or hardware.
\begin{itemize}
\item P and Q at 15 require 8 bytes at that cut. Adding E at 21 without
expiring anything needs 9 for both goals (the already due Q and correct
admission). A cap of 8 can admit Q and thereby crowd out the one-byte E; a
cap of 7 can instead reach correct admission after losing Q. Joint success is
not the same outcome as correctly admitting the inspection.
\item P at 5 still counts at 21; P at 4 no longer does. With Q at 15,
scheduling P at 4 rather than 5 permits both duties at cap 8. This relies on
the historical record and its receiver being available then. It cannot be
applied by moving the not-yet-settled outcome back in time.
\item With P and Q at 15 and cap 9, a bit-identical packet from another
sender W at 21 can consume the last byte before E arrives. The app correctly
rejects W, but that does not refund the already charged receive resource.
E-first admits the two duties; W-first needs one additional byte for both.
Admission at 21 is a declared cut, not a computed lower bound on whether an
uncommitted partial turn could be made before a later response.
\end{itemize}
An independent integer-clock accounting oracle agrees with the native broker
over 512 cases, including the two expiry sides, missing P, physically late E
and same-byte wrong-source controls. Native reachability checks are run before
callback-based selection and are not inputs to that selector. These are
selected-source comparisons, with no outperforming claim against an equally
informed direct scheduler. The unsettled assumption is how to provide a timely,
authorized source-to-plant effect record and stable endpoint in an independently
selected mission, not only in this simulator.
\end{document}
